% Source: Weinberg GR, physical PDF pp. 281--282 (printed pp. 259--260).
% Coverage: Section 10.3, Energy and Momentum of Plane Waves; equations
%           (10.3.1)--(10.3.7).
% Figures/tables/footnotes: none.
% Status: source-reviewed and compile-clean.
% Uncertainties: none.
% MODERNIZATION: R^(N) is the target-convention curvature and hence the
% negative of Weinberg's R^(N). The signs in (10.3.1), (10.3.2), and
% (10.3.4) have been converted together, leaving the positive physical
% energy-momentum formulas (10.3.5)--(10.3.7) unchanged.

\subsection{Energy and Momentum of Plane Waves}\label{sec:10.3}

The physical significance of the plane-wave solution \eqref{eq:10.2.1} is
brought forward by calculating the energy and momentum it carries. According
to Eq.~\eqref{eq:7.6.4}, the energy-momentum tensor of gravitation is given
to order \(h^2\) by
\[
  t_{\mu\nu}\simeq\frac{1}{8\pi G}
  \left[
    \frac12h_{\mu\nu}\eta^{\lambda\rho}R^{(1)}_{\lambda\rho}
    -\frac12\eta_{\mu\nu}h^{\lambda\rho}R^{(1)}_{\lambda\rho}
    -R^{(2)}_{\mu\nu}
    +\frac12\eta_{\mu\nu}\eta^{\lambda\rho}R^{(2)}_{\lambda\rho}
  \right],
\]
where \(R^{(N)}_{\mu\nu}\) is the term in the Ricci tensor of order \(N\)
in \(h_{\mu\nu}\). The metric
\(g_{\mu\nu}=\eta_{\mu\nu}+h_{\mu\nu}\) satisfies the first-order Einstein
equations \(R^{(1)}_{\mu\nu}=0\), so we can drop these terms in
\(t_{\mu\nu}\) and use
\begin{equation}
  t_{\mu\nu}\simeq
  -\frac{1}{8\pi G}
  \left[
    R^{(2)}_{\mu\nu}
    -\frac12\eta_{\mu\nu}\eta^{\lambda\rho}R^{(2)}_{\lambda\rho}
  \right].
  \tag{10.3.1}\label{eq:10.3.1}
\end{equation}
(For the actual metric it is \(R_{\mu\nu}\), rather than
\(R^{(1)}_{\mu\nu}\), that vanishes, and \(t_{\mu\nu}\) arises only from
the first-order terms in Eq.~\eqref{eq:7.6.4}. Here it is
\(R^{(1)}_{\mu\nu}\), rather than \(R_{\mu\nu}\), that vanishes, because
\(g_{\mu\nu}=\eta_{\mu\nu}+h_{\mu\nu}\) satisfies the first-order Einstein
equations rather than the exact equations. The difference is only of order
\(h^3\).)

To calculate \(R^{(2)}_{\mu\nu}\) we must use Eq.~\eqref{eq:10.2.1} in
Eq.~\eqref{eq:7.6.15}; the result is extremely complicated, but can be
simplified if we average \(t_{\mu\nu}\) over a region of space and time much
larger than \(\lvert k\rvert^{-1}\). (This is the way the energy and
momentum of any wave are usually evaluated.) The averaging kills all terms
proportional to \(\exp(\pm2\ii k_\lambda x^\lambda)\), and we are left with
only the \(x^\mu\)-independent cross-terms:
\begin{align}
  \left\langle R^{(2)}_{\mu\nu}\right\rangle
  ={}&-\operatorname{Re}\Bigl\{
  e^{\lambda\rho*}
  \bigl[
    k_\mu k_\nu e_{\lambda\rho}
    -k_\mu k_\lambda e_{\nu\rho}
    -k_\nu k_\rho e_{\mu\lambda}
    +k_\lambda k_\rho e_{\mu\nu}
  \bigr]
  \notag\\
  &+
  \left[
    e^\lambda{}_\rho k_\lambda
    -\frac12e^\lambda{}_\lambda k_\rho
  \right]^*
  \left[
    k_\mu e^\rho{}_\nu+k_\nu e^\rho{}_\mu-k^\rho e_{\mu\nu}
  \right]
  \notag\\
  &-\frac12
  \left[
    k_\lambda e_{\rho\nu}
    +k_\nu e_{\rho\lambda}
    -k_\rho e_{\lambda\nu}
  \right]^*
  \left[
    k^\lambda e^\rho{}_\mu
    +k_\mu e^{\rho\lambda}
    -k^\rho e^\lambda{}_\mu
  \right]
  \Bigr\}.
  \tag{10.3.2}\label{eq:10.3.2}
\end{align}
(We have not yet made use of the conditions \eqref{eq:10.2.2} and
\eqref{eq:10.2.3} appropriate to harmonic coordinates, so suppose for a
moment that we leave the harmonic coordinate system by adding to
\(h_{\mu\nu}(x)\) a term
\begin{equation}
  \ii(q_\mu\zeta_\nu+q_\nu\zeta_\mu)
  \exp(\ii q_\lambda x^\lambda)
  -\ii(q_\mu\zeta_\nu^*+q_\nu\zeta_\mu^*)
  \exp(-\ii q_\lambda x^\lambda),
  \tag{10.3.3}\label{eq:10.3.3}
\end{equation}
where \(q_\mu q^\mu\ne0\). After averaging over spacetime distances large
compared with \(\lvert q-k\rvert^{-1}\), the interference between
\eqref{eq:10.2.1} and \eqref{eq:10.3.3} drops out, and we find for
\(\langle R^{(2)}_{\mu\nu}\rangle\) the term \eqref{eq:10.3.2}, plus another
term obtained by replacing \(k\) with \(q\) and \(e_{\mu\nu}\) with
\(q_\mu\zeta_\nu+q_\nu\zeta_\mu\). Inspection of \eqref{eq:10.3.2} shows
immediately that this second term vanishes, so
\(\langle R^{(2)}_{\mu\nu}\rangle\) and hence
\(\langle t_{\mu\nu}\rangle\) may be calculated in harmonic coordinates with
no loss of generality.)

If we now use in \eqref{eq:10.3.2} the harmonic coordinate conditions
\eqref{eq:10.2.2} and \eqref{eq:10.2.3}, we find
\begin{equation}
  \left\langle R^{(2)}_{\mu\nu}\right\rangle
  =
  -\frac{k_\mu k_\nu}{2}
  \left(
    e^{\lambda\rho*}e_{\lambda\rho}
    -\frac12\left|e^\lambda{}_\lambda\right|^2
  \right).
  \tag{10.3.4}\label{eq:10.3.4}
\end{equation}
The quantity
\(\eta^{\lambda\rho}\langle R^{(2)}_{\lambda\rho}\rangle\) vanishes because
\(k^\rho k_\rho=0\), so \eqref{eq:10.3.1} now gives the average
energy-momentum tensor of a plane wave as
\begin{equation}
  \left\langle t_{\mu\nu}\right\rangle
  =
  \frac{k_\mu k_\nu}{16\pi G}
  \left(
    e^{\lambda\rho*}e_{\lambda\rho}
    -\frac12\left|e^\lambda{}_\lambda\right|^2
  \right).
  \tag{10.3.5}\label{eq:10.3.5}
\end{equation}
Note that a ``gauge transformation'' \eqref{eq:10.2.7} will change the
terms in \(\langle t_{\mu\nu}\rangle\) into
\[
  e'^{\lambda\rho*}e'_{\lambda\rho}
  =
  e^{\lambda\rho*}e_{\lambda\rho}
  +2\operatorname{Re}
    \bigl(\zeta_\rho^*k^\rho e^\lambda{}_\lambda\bigr)
  +2\left|\zeta_\rho k^\rho\right|^2,
\]
\[
  e'^\lambda{}_\lambda
  =
  e^\lambda{}_\lambda+2k^\lambda\zeta_\lambda,
\]
but \(\langle t_{\mu\nu}\rangle\) is gauge-invariant! Thus, as far as energy
and momentum are concerned, the polarizations \(e_{\mu\nu}\) and
\(e_{\mu\nu}+k_\mu\zeta_\nu+k_\nu\zeta_\mu\) represent the same physical
wave, and we see again that there are not six but only two physically
significant polarization parameters.

In particular, a wave traveling in the \(z\)-direction, with wave vector and
polarization tensor given by \eqref{eq:10.2.8} and \eqref{eq:10.2.9}, has
the energy-momentum tensor
\begin{equation}
  \left\langle t_{\mu\nu}\right\rangle
  =
  \frac{k_\mu k_\nu}{8\pi G}
  \left(\left|e_{11}\right|^2+\left|e_{12}\right|^2\right),
  \tag{10.3.6}\label{eq:10.3.6}
\end{equation}
or, in terms of the helicity amplitudes \eqref{eq:10.2.15},
\begin{equation}
  \left\langle t_{\mu\nu}\right\rangle
  =
  \frac{k_\mu k_\nu}{16\pi G}
  \left(\left|e_+\right|^2+\left|e_-\right|^2\right).
  \tag{10.3.7}\label{eq:10.3.7}
\end{equation}
